\section{问题分析}

\subsection{总体思路}

四个问题共享同一套微网物理结构。分别为每一问重新建模会重复书写能量平衡、储能容量和功率约束，也容易使效率定义、时间索引和计费方式前后不一致。本文先固定“电如何在微网中流动”这一物理内核，再随着题目推进改变“决策时能看到什么”和“计划改变后怎样结算”，并把微网能量平衡与储能状态递推作为公共模型，把预测、风险修正、滚动更新和价格变化分别放到对应问题中处理。

在这条主线上，四问形成清晰的递进关系。问题一先回答确定信息下储能与外网购电如何协同；问题二引入计划与实际之间的偏差，重点转向预测误差和高价紧急购电风险；问题三给予日内更新预报，使模型能够比较“继续执行旧计划”和“付出调整代价修改计划”之间的得失；问题四则检验同样的控制逻辑在逐日电价变化后是否仍成立。全文据此按照“统一数据基准—建立公共物理模型—逐问扩展信息和费用规则—解释调度机理—验证结果可靠性”的顺序展开。

\subsection{各问题的建模重点}

\indent\textbf{问题一：}附件1给出了典型日的电价、负载和光伏预测，可直接建立确定性经济调度模型。此时储能的收益来源十分明确：一是把低价时段购买的电转移到高价时段使用，二是把中午暂时超过负载的光伏保存到之后。题目又要求 $0{:}00$ 与 $24{:}00$ 储电量相同，因此不能通过消耗初始库存人为降低费用；首末储电状态必须闭合。

\indent\textbf{问题二：}本问与问题一的关键差别在于全天计划必须在源荷真实实现之前确定。主模型严格遵循因果性，只利用 $0{:}00$ 之前已经发生的数据形成当天负载和光伏预测；若实际业务中能够提前获得全天计划负荷，则把该情形作为信息更充分的对照，不与主模型混用。紧急购电价格达到正常价格的 $5$ 倍，高估光伏会比同量低估产生更大的经济损失，预测也就不能只追求平均误差最小，还需要体现这种非对称风险。全天计划确定后不再修改外网购电量，但储能可根据当前已实现偏差调整当期充放电，使“计划层决策”和“执行层补偿”保持清晰分工。

\indent\textbf{问题三：}本问的核心是判断新预测到达后，重新调整未来购电计划是否划算。每天 $0{:}00$ 先形成覆盖全天的基准计划，$6{:}00$、$12{:}00$、$18{:}00$ 到来时只重算尚未执行的区间，并从当前实测储电量继续递推。为了回答题目关于“其他时刻预报是否需要”的问题，必须保持物理约束、预测器和结算规则不变，只逐步增加允许使用的更新节点，再比较全年费用、紧急购电和调整量。由于节点到来时不仅有新光伏预报，也多了已经实现的负载和状态信息，节点收益应解释为综合决策信息的价值，并进一步用配对对照区分各信息来源。

\indent\textbf{问题四：}题面已给出附件4的全年实时电价，正式重算直接以其替换前几问的典型日电价，其他源荷预测、储能反馈和滚动逻辑保持不变。这样，前后结果的差异可以主要归因于价格环境改变。考虑真实市场中未来价格未必在日初全部公开，再设置仅用历史价格预测未来价格的扩展情形，用于衡量价格信息受限的机会成本，而不替代题面主结果。

\subsection{求解方法选择}

上述能量平衡、状态递推、容量边界以及分段线性费用都可以写成线性约束。问题一、问题二及不要求充放电严格互斥的对照方案属于线性规划；当需要从物理上禁止同一时段发生充、放电重叠时，引入二元变量后形成混合整数线性规划。两类模型均由 HiGHS 求解。相较于遗传算法、粒子群等元启发式方法，线性/混合整数规划可以直接利用本题的结构，在给定规则下获得可验证的最优解，并能逐时回代检查每一项能量和费用，这与本文强调的可解释性和可复现性更一致。
